\chapter{Slutsats}

Social engineering har i denna avhandling analyserats som ett cybersäkerhetshot med fokus på samverkan mellan psykologiska och tekniska aspekter av digital manipulation.
Den första forskningsfrågan har besvarats genom att analysera hur social engineering bygger på systematiska kognitiva biaser och sociala principer såsom auktoritet, brådska,
tillit och social bekräftelse. Dessa psykologiska mekanismer påverkar individers beslutsfattande och minskar benägenheten att kritiskt granska situationer som framstår som legitima.
Angreppens effektivitet ligger därmed inte endast i teknisk komplexitet, utan också i förmågan att manipulera mänskliga beteendemönster och organisatoriska normer.

Den andra forskningsfrågan har besvarats med att visa hur moderna social engineering-attacker ofta kombinerar traditionella metoder såsom nätfiske, spjutfiske, pretexting, baiting 
och insider-relaterade attacker med tekniska komponenter som artificiell intelligens, domänförfalskning och automatiserade meddelandesystem. Fallstudierna illustrerar även hur
AI-baserade djupfejkar och riktade informationsoperationer ytterligare förstärker attackernas trovärdighet och effektivitet. Social engineering-attacker visar därmed hur
flexibla och anpassningsbara de är, då både tekniska och sociala element kombineras för en unik angreppsstrategi.

Den tredje forskningsfrågan har besvarats genom att analysera olika detektionssystem som främst bygger på tekniska filtreringssystem, maskininlärningsbaserad analys, SIEM- och IDS-integration
samt automatiserade responsmekanismer såsom SOAR. Dessa system kan upptäcka avvikande beteenden, misstänkt kommunikation och ovanliga åtkomstförsök. Det visades också att systemen
måste kompletteras med organisatoriska rutiner och säkerhetskultur, eftersom manipulationen ofta riktas mot mänskligt beslutsfattande snarare än enbart tekniska sårbarheter.

Den fjärde forskningsfrågan har besvarats genom att analysera den artificiella intelligensens potential inom detektion och respons genom avancerad mönsterigenkänning och beteendeanalys.
AI kan bidra till snabbare identifiering av cyberhot och effektivare hantering av incidenter. Däremot begränsas samtidigt teknologins effektivitet av faktorer som datakvalitet, falska resultat och
det faktum att hotaktörer också använder AI för att förbättra sina angreppsstrategier. AI bör därför förstås som en relevant del av försvarsstrategier mot social engineering, men inte
som en självständig lösning.
